% ==================================================================
%---------BLOQUE 1------------------- (Preambulo y datos del documento)
%====================================================================
% NOTA PARA EL FORO (UNA PÁGINA), CA-0309 CONTINGENCIAS DE VIDA II
% Grupo 2: Auditoría o Control, ronda 1
%
% Compilar con LuaLaTeX, dos veces.
% La tabla la genera Codigo/main.R en salidas/tablas
%====================================================================



% !TEX program = lualatex
\documentclass[10pt,a4paper]{article}

% ---------- Matemática y fuentes ----------
\usepackage{amsmath,mathtools}
\usepackage{fontspec}
\usepackage{unicode-math}
\setmathfontface\mathsf{Latin Modern Sans}   % Sans-serif matemático
\setmathfontface\mathtt{Latin Modern Mono}   % Typewriter matemático
\setmathfont{NewCMMath-Book}
% ---------- Idioma ----------
\usepackage{polyglossia}
\setmainlanguage{spanish}
\addto\captionsspanish{%
	\renewcommand{\tablename}{Tabla}%
}
% ---------- Diseño ----------
\usepackage[a4paper,top=1.7cm,bottom=1.7cm,left=1.4cm,right=1.4cm,headsep=0.35cm]{geometry}
\setlength{\emergencystretch}{2.5em} % evita líneas que se salen
\usepackage{microtype}
\usepackage{enumitem}
\usepackage[table]{xcolor}
\usepackage{booktabs}
\usepackage{fancyhdr}
\setlength{\headheight}{11pt}
\pagestyle{fancy}
\fancyhf{}
\renewcommand{\headrulewidth}{0pt}
\fancyhead[R]{\small\itshape Nota para el foro}
\fancyfoot[C]{\small\thepage}

\setlist[itemize]{noitemsep, topsep=1pt, leftmargin=1.2em}
\setlist[enumerate]{noitemsep, topsep=1pt, leftmargin=1.4em}

% ---------- Enlaces ----------
\usepackage[colorlinks=true,linkcolor=blue!40!black,citecolor=blue!40!black,urlcolor=blue!40!black]{hyperref}

% ---------- Macros ----------
% Tablas exportadas por R (solo el tabular)
\newcommand{\tabla}[1]{\input{salidas/tablas/#1.tex}}

% ---------- Encabezado a una sola columna (franja centrada) ----------
% Misma franja del informe, sin resumen ni palabras clave.
\newcommand{\anchofranja}{0.78\textwidth}
\makeatletter
\newcommand{\afiliaciones}[1]{\gdef\@afiliaciones{#1}}
\let\@afiliaciones\@empty
\renewcommand{\@maketitle}{%
	\newpage\null\vskip -1.5em
	\begin{center}
		\begin{minipage}{\anchofranja}
			\centering
			{\large \@title\par}
			\vskip 0.6em
			{\def\and{\unskip,\hskip 0.8em plus 0.3em\ignorespaces}\small \@author\par}
			\ifx\@afiliaciones\@empty\else
			\vskip 0.3em
			{\footnotesize \@afiliaciones\par}
			\fi
			\vskip 0.4em
			{\footnotesize \@date\par}
		\end{minipage}
	\end{center}
	\vskip 0.3em}
\makeatother



%====================================================================
% DATOS DEL DOCUMENTO
%====================================================================
\title{\textbf{Nota para el foro: robustez de la propuesta ante mayor renuncia} \\[0.2em]
	\normalsize Grupo 2, Auditoría o Control \,|\, Ronda 1}

\author{%
	Andrey Gonzales (C33329)\textsuperscript{1} \and
	Randal Picado Bermúdez (C36024)\textsuperscript{1} \and
	Anthonny Flores Rojas (C32975)\textsuperscript{1}}

\afiliaciones{%
	\textsuperscript{1}Escuela de Matemática, Universidad de Costa Rica}

\date{Curso: Contingencias de Vida II (CA-0309) \,|\, Profesora: Yerling Ramírez Morera \,|\, \today}



%---------BLOQUE 2------------------- (Contenido de la nota)
\begin{document}

	\maketitle
	\thispagestyle{fancy}

	\noindent\textbf{Escenario.} Prima anual de 55\,000 y renuncia de decremento único del 6\,\%, con los demás parámetros de la base (rendimiento 5\,\%, gasto 10\,000, descuento 4\,\% y el mismo cuadro de reservas por contrato). Se recalcularon probabilidades competitivas, población, flujos, reservas totales y recursos. El 6\,\% es un estrés académico, no una reacción estimada de los afiliados. Montos en millones de u.m.

	\begin{table}[h]
		\centering\small
		\tabla{S4_base_grupo2}
		\caption{Caso base frente al escenario del grupo 2 (prima de 55\,000 y renuncia del 6\,\%).}\label{tab:grupo2}
	\end{table}
	\vspace{-0.4em}

	\noindent\textbf{Comparación con nuestra base.} El efecto conjunto mejora el NPV en unos 2\,690 millones y elimina los aportes, pero el producto sigue perdiendo (margen \(-44.09\,\%\)) y los recursos siguen sin respaldar la reserva al cierre del año 1. La cartera termina con 12\,312 activos menos. Hay unos 373 fallecimientos cubiertos menos porque hay menos asegurados, no porque el riesgo por contrato baje: el siniestro esperado por contrato vigente (45\,900 a 56\,500) sigue por encima de la nueva prima neta de 42\,250. No atribuimos la mejora a la prima o a la renuncia por separado.

	\medskip
	\noindent\textbf{Observación 1: la prima no cubre el siniestro esperado.}
	\begin{itemize}
		\item \textit{Evidencia:} en la base, la prima neta de gasto y comisión es 37\,500 por contrato y el siniestro esperado va de 46\,659 a 57\,392. Aun sin adquisición, la operación pierde 5\,209.95 en valor presente; el NPV es \(-12\,912.86\) (margen \(-57.88\,\%\)) y llega a \(-13\,617.35\) en el adverso.
		\item \textit{Consecuencia:} ningún escenario analizado es rentable; la pérdida es estructural y crece con la edad de la cartera.
		\item \textit{Acción:} re-tarificar con una prima que cubra siniestro y gasto por edad, y revisar la adquisición de 75\,000, que explica el 60\,\% de la pérdida base.
	\end{itemize}

	\medskip
	\noindent\textbf{Observación 2: los recursos no respaldan la reserva desde el año 1.}
	\begin{itemize}
		\item \textit{Evidencia:} \(\widetilde F_1-V_1\) es \(-2\,024.61\) en la base, \(-2\,197.78\) en el adverso y \(-1\,307.79\) en nuestro escenario. La base requiere 3\,050.20 de aportes y el adverso 3\,807.06, a partir del año 3.
		\item \textit{Consecuencia:} la aseguradora dependería de capital externo para pagar siniestros y no tendría activos que respalden la reserva.
		\item \textit{Acción:} exigir un plan de capital inicial aprobado o reaseguro antes del lanzamiento, y verificar \(\widetilde F_t\ge V_t\) en cada cierre.
	\end{itemize}

	\medskip
	\noindent\textbf{Nuestras preguntas.}
	\begin{enumerate}
		\item Al grupo actuarial: ¿qué evidencia sostiene la renuncia del 3\,\%, y cómo cambia su recomendación si un alza de prima provoca más renuncias y se quedan los afiliados con mayor mortalidad?
		\item Al grupo actuarial y al Comité: ¿con qué capital se cubriría el déficit de \(-2\,024.61\) millones frente a la reserva en el año 1 y los 3\,050.20 millones de aportes, y qué prima mínima consideran para que el NPV no sea negativo?
	\end{enumerate}

	\medskip
	\begin{center}
		\fbox{\begin{minipage}{0.96\textwidth}
				\textbf{Opinión preliminar: posponer} el seguro colectivo. Con la tarifa actual pierde en base, adverso y estrés, y no respalda la reserva desde el año 1. Lo apoyaríamos con NPV no negativo en base y adverso, un plan de capital para el respaldo y seguimiento de la renuncia para no ganar a costa de dejar afiliados sin protección. \textbf{Apoyar} la cuenta equity-linked con \(G=100\,000\) (NPV medio 4\,054.40 frente a 1\,093.45; CTE\(_{95}\) 22\,959.79 frente a 29\,892.45), condicionada a monitorear la cola. \textit{Límites:} supuestos de la guía, reservas suministradas, población cerrada y una sola muestra de 1\,000 trayectorias.
		\end{minipage}}
	\end{center}

\end{document}
